\scanpage{109}
\mspage{54}

(C'est un foncteur p.r. au triple $\widetilde X\longrightarrow X\longrightarrow Y$ avec $X,Y$ $0$-\unsure{connexes|connexe}, $\widetilde X$ un rev. universel; plus généralt, si $\deleted{Z}\inserted{T}\longrightarrow X$ avec $\deleted{Z}\inserted{T}$ $1$-connexe, on pose
\[
(11)\qquad \unsure{\mathcal E}_{X/Y}^{\deleted{Z}\inserted{T}}=\operatorname{Ker}\bigl(\mathcal E_X^{\deleted{Z}\inserted{T}}\longrightarrow\mathcal E_Y^{\deleted{Z}\inserted{T}}\bigr)
\]
on a un foncteur en \inserted{le triple} $(\deleted{Z}\inserted{T}\longrightarrow X_{\unsure{0|\bullet}}\longrightarrow Y)$ variable. En particulier: $\mathcal E_{X/Y}^{\unsure{\zeta|\xi}}$, $\unsure{\zeta|\xi}$ un $p^t$ géom. de $X$, $\mathcal E_{X/Y}^{\Omega}$, $\mathcal E_{X/A}^{\widetilde X}$ (si $Y=\operatorname{Spec} A$), $\mathcal E_{B/A}^{\widetilde B}$ $\ldots$

\margin{\uline{NB} $\mathcal E_{X/Y}^T$, pour $T$ variable, $X/Y$ fixé, est défini à isom non unique !) près, et se note $\mathcal E_{X/Y}$ \unsure{simplem\textsuperscript{t}.}}

Un cas intéressant est celui où on dispose d'une ``suite exacte d'homotopie universelle'' \deleted{de} (en dim $1$) pour $X\longrightarrow Y$, alors la donnée (pour $X=Y$ donné) de $T\longrightarrow X$, (avec $T$ $1$-connexe) peut s'interpréter par la donnée du composé

\margin{Car \unsure{où|à} $\mathcal E_X^{\deleted{Z}\inserted{T}}\longrightarrow\mathcal E_Y^{\deleted{Z}\inserted{T}}$ est \unsure{épi|épi.}}
\margin{108}
